% lec6.tex 
\lecture{6}{08:05 AM Thu, Nov 06 2025}{} 
\begin{theorem}[Chinese remainder theorem]
  Let $p,q \in  \NN, (p, q) = 1$. Then $\mathbb{Z}_{pq}$ is isomorphic to $\mathbb{Z}_{p} \times \mathbb{Z}_{q} $.
\end{theorem}
\begin{proof}
Let \[
\begin{array}{cccc}
      f : &  \mathbb{Z}  & \longrightarrow & \mathbb{Z}_{p} \times 
      \mathbb{Z}q\\

           &  k  & \longmapsto     & (\overline{k}^{p}, \overline{k}^{q})  \\ 
\end{array}
\]
We have $f$ is a RM, also,
\begin{align*}
\text{Ker}  (f) 
&= 
\left\{ k \in   \mathbb{Z}: \overline{k}^{p} = \overline{0}^{p}
\text{ and }  \overline{k}^{q}= \overline{0}^{q}\right\} \\
&= 
\left\{ k \in   \ZZ: 
k \in   p\ZZ \cap q \ZZ\right\} \\
&= \text{lcm} (p , q)\ZZ = pq \ZZ.
\end{align*}
Hence the kernel of $f$ is $\left\{ \overline{0}^{pq} \right\}$. Then 
by first isomorphism theorem, 
\[
\ZZ_{pq} \text{ is isomorphic to } \Img (f)  \subset 
\ZZ_{p} \times \ZZ_{q}.
\]
so $pq = \left| pq \ZZ \right|  = \left| \Img (f)  \right|  =
\left| \ZZ_{p} \times \ZZ_{q}  \right|  $. Then 
$\ZZ_{pq}$ is isomorphic to $\Img (f) \times \ZZ_{p} \times \ZZ_{q}.$   
\end{proof}
\begin{proposition}[]
Let $C_{p}$ and $C_{q}$ be two cyclic groups. Then 
$C_{p} \times C_{q} $ is cyclic if and only if 
$(p, q) = 1.$ 
\end{proposition}
\begin{proof}
  \ding{50} $( \impliedby  ) $: $(p, q)=1 \iff C_{p} \times C_{q} $ is 
  isomorphic to $Z_{p} \times Z_{q}$ which is isomorphic to 
  $Z_{pq}$. Thus, its cyclic. \\\ding{50}$( \implies ) $: 
  $d = (p, q)  > 1$. Then $\text{lcm}  (p, q) = a < pq$. Let
  $(\overline{x}^{p}, \overline{x}^{q}) \in  \CC _{p} \times C_{q} $ 
  which is isomorphic to $Z_{p}\times \ZZ_{q}$. Then 
  $(\overline{x}^{p}, \overline{y}^{q}) = (\overline{ax}^{p}, 
  \overline{ay}^{q}) = (\overline{0}^{p}, \overline{0}^{q}). $ So
  $C_{p} \times C_{q} $ is not cyclic.
\end{proof}
\divider
\subsection{Euler function}
\noindent\ding{43}~\textbf{Euler function:} for $n \in  \NN$, define the function $\FF$ by:
\[
\begin{array}{cccl}
      \FF : &  \NN  & \longrightarrow & \NN \\

           &  n  & \longmapsto     & \FF(n)  = 
           \left| \left\{ 1 \leq k \leq n: (k, n)  = 1 \right\} \right|  \\ 
\end{array}
\]
\ding{100}~\textbf{Properties. } If $p$ is prime, $\FF(p) = p -1$. 
for all $r \in  \NN$, we have, 
\[
\FF( p ^r ) = 
p ^{r} - p ^{r-1} = p ^{r-1} ( p - 1).
\]
If $n,m \in  \NN$ such that $(n, m) = 1$, then
\[
\FF(nm) = \FF(n)  \FF(m). 
\]
\begin{proof}
  \begin{align*}
    \FF(p ^r)  &= 
  \left| \left\{ 1 \leq k \leq p ^{r} : (k, p ^{r}) = 1 \iff 
  (k, p) = 1\right\} \right|  \\
               &= 
               p ^{r} - 
               \left| \left\{ 1 \leq k \leq  p ^{r} : 
               (k, p) = p\right\} \right|  \\
               &= 
               p ^{r} - p ^{r-1}
  \end{align*}
  We can use the chinese remainder theorem, 
  \begin{align*}
    (n, m)  & \implies  1 \implies 
    \ZZ_{nm}\text{ is isomorphic }  
    \ZZ_{n} \times \ZZ_{m}\\
            & \implies  
            (\ZZ_{nm})^{*} 
            \text{ is isomorphic }
            (\ZZ_{n} \times \ZZ_{m} ) ^{*} \\
            & = \ZZ_{n}^{*}\times \ZZ_{m}^{*} 
  \end{align*}
  Alternative way, we have $\FF(nm) = \left| (\ZZ_{nm}) ^{*} \right| = 
  \left| \ZZ_{n}^{*} \times \ZZ_{m} ^{*}  \right| = 
  \left| \ZZ_{n}^{*} \right| \cdot  
  \left| \ZZ_{m}^{* } \right| = \FF(n) \times \FF(m)$. 
\end{proof}
\noindent Let $n \geq 2$ (where $n \in  \NN$).\\
\\
\ding{95}~\textbf{Proposition. }
\begin{itemize}
  \item[\ding{172} ] 
    Let $a \in  \ZZ$, if $(a, n) = 1$ then $a ^{\FF(n) } = 1
    \left[ n \right]$.
  \item[\ding{173} ] 
    Every cyclic group of order $n$ has $\FF(n) $ generators.
\end{itemize}
\begin{proof}
  \ding{50} \ding{172}. Let $a \in \ZZ$, and suppose that $(a, n) = 1$, that is; 
\begin{align*}
  (a, n)  = 1 
  & \implies 
  \exists \al, \bb \in  \ZZ: a \al + \bb n = 1, \\
  & \implies 
  \overline{\al} \overline{a} = \overline{1} \in  
  \ZZ_{n}, 
  \\
  & \implies 
  \overline{a} \in  
  \left( \ZZ_{n} \right) ^{*},
  \\
  & \implies 
  \overline{a}^{\FF(n) } = 
  \overline{1} \quad \left( \left| \ZZ_{n}^{*} \right| =
  \FF(n) \right), \\
  & \implies 
  a ^{\FF(n) } = 
  1 \left[  n \right].
\end{align*}
\ding{50} \ding{173}. Left as an exercise \textbf{\ding{106}}. This finishes
the proof, as required.
\end{proof}
\divider
We denote the finite field $\ZZ_{p} := \mathbb{F}_{p}.$ \\
\begin{proposition}[]
Let $n \geq 2$ (where $
n \in  \NN$). Then $\ZZ_{n}$ is a field if and only if 
$n$ is a prime.
\end{proposition}
\begin{proof}
If $n$ is prime, then $\left| \ZZ_{n} \right|  ^{*} = 
\FF(n)  = n-1.$ So $\ZZ_n$ is a field. Suppose $\ZZ_n $ is a 
field. Let $1 \leq  d < n$ such that $d|n$. So there exists 
$q \in  \NN$ such that $ n = q d.$ Then $\overline{q} \overline{d}= 
\overline{0} \in  \ZZ_{n}$; Since $d \neq 0$, then 
$\overline{d} \in  (\ZZ_n ) ^{*}.$  So $\overline{q} = \overline{0}$, thus
$q = n$ and $d= 1$. Hence $n$ is prime.
\end{proof}
\begin{theorem}[Fermat's Theorem]
  Let $p$ be prime. Let also $a \in  \ZZ$, then 
  \[
    a ^{p} = a \left[ p \right].
  \]
\end{theorem}
\begin{proof}
  Since $\ZZ_{p}$ is a field, then for all $a \in  \ZZ \backslash  p \ZZ$, 
  we have:
\end{proof}
\[
\overline{a}^{p-1} = \overline{1} \in  
\left( \ZZ_{p} \right) ^{*} \implies 
\overline{a}^{p} = \overline{a}.
\]
If $a \in  p\ZZ$, then $\overline{a}^{p} = \overline{0}^{p} = \overline{0}
= \overline{a}$. Thus
\[
  a ^{p} = a \left[ p \right], \quad ( 
  \forall a \in  \ZZ).
\]
\begin{definition}[]
Let $R$ be a ring. We say that $R$ is an integral domain 
(ID) if:
\begin{itemize}
  \item[\ding{172} ] $R$ is
    commutative.
  \item[\ding{173} ] $R$ is
    unity.
  \item[\ding{174} ] $R$ has 
    no zero divisors; That is, for $a, b \in  R$:
    \[
    ab = 0 \implies a = 0 \text{ or }  b = 0.
    \]
\end{itemize}
\end{definition}
\section{Principal ring}
\begin{definition}[]
A ring $R$ is a principal ring if $R$ is an integral domain, and
every deal of $R$ is principal. That is; for all $I \subset R$, 
if $I$ is an ideal, then:
\[
I = \left( a \right)  = a R, \quad (a \in  R).
\]
\end{definition}
\begin{example}
  \begin{itemize}
    \item [\ding{172} ] $\ZZ$ is a principal ring 
      since $n \ZZ = (n)$ (where $n \in  \NN$), which are
      the ideals of $\ZZ.$ 
    \item [\ding{173} ] If $\KK$ is a field, then 
      $\KK \left[ X \right]$ is a principal ring. 
      Hint: (Use 
      the Euclidean division for the proof).\normalfont.
    \item [\ding{174} ] 
      The ring $\ZZ\left[ X \right]$ is not a principal ideal 
      since $(2, X) $ is not a principal ideal of $\ZZ \left[ X \right]$. 
      $\left( (a,b) = a R + b R \right) $.
  \end{itemize}
\end{example}
\begin{definition}[]
Let $a, b \in  R$. Suppose that $a = b c$ for some $c \in  R^{*}$,
then we say that $a$ and $b$ are associates.
\end{definition}
\noindent\textbf{Remark. }The relation ($a$ and $b$ are associates where $a, b \in  R$)
is an equivalence relation on $R$.
\begin{proposition}[]
  Let $R$ be a ring. Suppose that $R$ is an integral domain, then
  \[
    a R \subset b R \iff b \backslash  a, \quad 
    (a, b \in  R),
  \]
  and
  \[
    a R = b R \iff a = b c, \quad (c \in  R).
  \]
\end{proposition}
\subsection{Irreducible and prime elements}
\begin{definition}[]
Let $R$ be a ring. Suppose that $R$ is an integral domain, an element
$a \in  R \backslash  \left\{ 0 \right\}$ is said to be 
irreducible if $ a \notin  R^{*}$ (not invertible), and 
\[
a = uv \implies u \in  R ^{*} \text{ or }  v \in  R^{*},
\quad (u, v \in  R).
\]
\end{definition}
\begin{example}
\begin{itemize}
  \item The primes of $\ZZ$ are irreducible. Let $p$ be a prime number, and 
    let also $a, b \in  \ZZ$ such that
    \begin{align*}
    p = ab &\implies  a = p \text{ or }  b = p \\
           & \implies 
           b = 1 \in  \ZZ^{*} \text{ or }  
           a = 1 \in  \ZZ^{*}
    \end{align*}
   \item Let $\KK$ be a field. A polynomial $P \in  \KK\left[ X \right]$  of degree $n$ (where $n \in  \NN_0$) is irreducible if and only if 
     $n \geq 1$ and the only divisions are those of degree 
     $0$ or $n$. 
    \item In $\ZZ\left[ i \right]$, we have 
      $\left( 3 + 2i \right) \left( 3 - 2i \right) = 13$. 
      $\left( \ZZ\left[ i \right]^{*} = \left\{ \pm 1, \pm i \right\}  \right) $.
      So $13$ is not irreducible in $\ZZ\left[ i \right]$, even though 
      $13$ is irreducible in $\ZZ.$  
\end{itemize}
\end{example}
\begin{definition}[]
Let $R$ be a ring. Suppose that $R$ is an integral domain and 
let $p \in R$ such that $p \notin  R^{*} \cup \left\{ 0 \right\}$. If 
\[
p|ab \implies 
p|a \text{ or }  p|b, \quad  \left( a, b \in  R \right).
\]
Then we say that $p$ is prime. 
\end{definition}
% end of file
